\section{Conclusiones}

El objetivo general de este trabajo fue diseñar un esquema de control de posición para el sistema de levitación magnética (SLM) ECP-730 en configuración inestable considerando el fenómeno de atascamiento-deslizamiento. A continuación se resumen los resultados obtenidos para cada uno de los cuatro objetivos específicos, todos ellos obtenidos por simulación.

\paragraph{Análisis de la dinámica del SLM (primer objetivo).}
Se adoptó el modelo mecánico $m\ddot{x} = u/[b(a-x)^4] - m c_1\dot{x} - mg$, con la fuerza electromagnética del fabricante validada experimentalmente por Alcántara et al.\ \cite{alcantara} y los parámetros identificados en \cite{diana_tesis} (tabla~\ref{tab:parametros}). El análisis mostró que, en la configuración atractiva, el sistema admite un punto de equilibrio para cada posición y que todos ellos son inestables, puesto que $a - x_1^* > 0$ en todo el intervalo alcanzable. El efecto atascamiento-deslizamiento se modeló como la combinación de una zona muerta no centrada en el origen, cuyo centro se desplaza con el voltaje de equilibrio de cada posición, y de la fricción estática representada con el método de Karnopp. El jacobiano del simulador no lineal coincide con el modelo linealizado, lo que valida la implementación numérica. En lazo abierto, una perturbación de $0{,}01$~cm alrededor del equilibrio alcanza 1~cm en unos $0{,}3$~s (figura~\ref{fig:p8_lazo_abierto}); sin embargo, con el levitante en reposo, la fricción estática lo retiene en una banda de unos $0{,}8$~cm alrededor de cada equilibrio (figura~\ref{fig:p8_retrato}), de modo que la inestabilidad solo se manifiesta una vez que el imán se pone en movimiento.

\paragraph{Diseño del controlador sobre el modelo linealizado (segundo objetivo).}
El modelo linealizado en $x_1^*=-4$~cm tiene los polos $\lambda_1=-23{,}5057$ y $\lambda_2=14{,}9427$ y es completamente controlable y observable. Sobre este modelo se diseñó el controlador con doble acción integral
\[
C(s)=\frac{152{,}43\,(s+4{,}875)(s+1{,}861)(s+17{,}31)}{s^{2}(s+301{,}1)},
\]
que estabiliza el lazo con un margen de fase de $59{,}9^\circ$ en $129{,}3$~rad/s y un margen de ganancia inferior de $-17{,}1$~dB en $9{,}63$~rad/s, este último consecuencia del polo inestable de la planta (figuras~\ref{bode_fase} y~\ref{nyquist}). En lazo cerrado lineal, la respuesta al escalón presenta un sobrepaso de $16{,}9$\,\% y un tiempo de establecimiento de $0{,}284$~s (figura~\ref{step_response}). El mismo controlador mantiene la estabilidad en el intervalo de $-3$ a $-2$~cm, con márgenes de fase entre $54{,}5^\circ$ y $47{,}4^\circ$.

\paragraph{Restricción del actuador.}
La simulación con el modelo corregido mostró que el intervalo de la señal de control determina las posiciones en las que el sistema puede operar. Con $u\in[0;\,3{,}5]$~V, el levitante puede sostenerse en $x_1=-4$~cm, cuyo voltaje de equilibrio es $3{,}48$~V, pero no desplazarse desde ahí: para vencer la fricción estática hacia la bobina se requiere aproximadamente $1{,}145$ veces el voltaje de equilibrio, es decir, unos $3{,}99$~V. Mediante un barrido de 7503 simulaciones se determinó que los cambios de referencia son viables a partir de $x_1\approx-3{,}6$~cm; por ello, las pruebas de regulación y seguimiento se realizaron en el intervalo de $-3$ a $-2$~cm.

\paragraph{Simulación no lineal sin atascamiento-deslizamiento (tercer objetivo).}
Sobre el modelo no lineal sin fricción estática ni zona muerta, el PII lleva el levitante de $-2$ a $-3$~cm con un sobrepaso de $23{,}7$\,\% y entra en la banda de $\pm2$\,\% en $0{,}48$~s; el error converge a cero y el voltaje al valor de equilibrio de $2{,}39$~V, sin saturar (figuras~\ref{fig:sinad_pos} a~\ref{fig:sinad_ctrl}). El sobrepaso queda entre los que predice el modelo linealizado en $-3$~cm ($19{,}2$\,\%) y en $-2$~cm ($24{,}7$\,\%), lo que confirma la coherencia entre el modelo de diseño y la planta no lineal.

\paragraph{Simulación no lineal con atascamiento-deslizamiento (cuarto objetivo).}
Con zona muerta y fricción de Karnopp, el PII estabiliza el sistema inestable en las tres pruebas realizadas sin compensar explícitamente ninguna de las dos no linealidades. En regulación (figuras~\ref{posicion} a~\ref{senal_control}), el transitorio es prácticamente igual al del caso sin atascamiento-deslizamiento (sobrepaso de $23{,}6$\,\%), pero en estado estacionario la convergencia asintótica se sustituye por un ciclo límite de $0{,}033$~cm pico a pico, con el levitante atascado el 74\,\% del tiempo y un error medio prácticamente nulo. En seguimiento, el error se mantiene acotado en $\pm0{,}020$~cm para la referencia senoidal y en $\pm0{,}022$~cm para la trapezoidal, con errores eficaces de $0{,}012$ y $0{,}011$~cm. El ciclo límite está presente durante toda la prueba, con unas dos rupturas por segundo, y no solo cuando la referencia cambia de sentido. En todas las pruebas la señal de control se mantuvo por debajo de $3{,}3$~V, sin alcanzar el límite superior del actuador; solo en el instante del escalón cae momentáneamente al límite inferior de 0~V. La misma prueba de regulación sobre el modelo linealizado con zona muerta y fricción de Karnopp (figura~\ref{fig:p8_lineal_ad}) produce un transitorio y un ciclo límite casi idénticos ($0{,}028$ frente a $0{,}033$~cm), lo que indica que el ciclo límite se debe a la zona muerta y a la fricción estática, y no a la no linealidad de la fuerza electromagnética.

\paragraph{Configuración estable.}
En la configuración repulsiva, estable en lazo abierto con polos $-4{,}28\pm21j$, el mismo PII diseñado para la configuración inestable estabiliza el lazo con un margen de fase de $43{,}3^\circ$ y lleva el levitante de 1 a 2~cm con un sobrepaso de $9{,}8$\,\% en presencia de atascamiento-deslizamiento (tabla~\ref{tab:p8_estable}); también aquí el estado estacionario presenta un ciclo límite, de $0{,}020$~cm pico a pico.

\paragraph{Papel de la doble acción integral.}
La comparación con un controlador PI en las mismas condiciones (tabla~\ref{tab:pi_pii}) no confirmó las ventajas que motivaron el uso de la doble acción integral. El PII reduce ligeramente el sobrepaso ($23{,}6$\,\% frente a $25{,}9$\,\%), pero la duración media de cada atascamiento es mayor ($0{,}36$~s frente a $0{,}26$~s en regulación) y la amplitud del ciclo límite es prácticamente la misma. Un barrido de 256 variantes estables de ambos controladores (figura~\ref{fig:r15_familias}) muestra que no se trata de un efecto de la sintonización: ningún PII alcanza un ciclo límite menor que el mejor PI ($0{,}026$ frente a $0{,}013$~cm), y la amplitud del ciclo límite depende sobre todo de la ganancia proporcional. La eliminación del error en rampa que garantiza la segunda integración sí se observa, pero incluso con una pendiente doce veces mayor que la de las pruebas de seguimiento el error medio del PI ($0{,}56\times10^{-3}$~cm) es más de diez veces menor que el error eficaz que impone el ciclo límite. El análisis de la sección~\ref{sec:analisis_ciclo} explica el mecanismo: la fricción estática equivale a una banda muerta de voltaje de $\pm14{,}5$\,\% alrededor del equilibrio, la acción integral es la que obliga al imán a romperla una y otra vez (un PD con prealimentación no presenta ciclo límite), la duración de cada atascamiento la fija el término integral simple y la amplitud del ciclo es aproximadamente inversamente proporcional a la ganancia del controlador.

\paragraph{Relación con los antecedentes.}
Alcántara et al.\ \cite{alcantara} observaron en la configuración repulsiva del ECP-730, estable en lazo abierto, que la acción integral del lazo externo interactúa con la fricción y produce oscilaciones de tipo ciclo límite, y las eliminaron con control conmutado. Este trabajo encuentra el mismo mecanismo en la configuración atractiva, inestable en lazo abierto, y muestra que un controlador lineal con acción integral y sin conmutación basta para estabilizar el sistema en presencia de zona muerta y fricción estática, a costa de ese ciclo límite. En cuanto a la hipótesis de Pérez-Gómez \cite{perez}, según la cual los controladores con doble acción integral reducen la amplitud de las oscilaciones producidas por la zona muerta, los resultados obtenidos no la confirman para el SLM: en ninguna de las variantes evaluadas la doble acción integral reduce la amplitud del ciclo límite respecto a un PI.

\paragraph{Evaluación de la hipótesis.}
La primera parte de la hipótesis se confirma: el PII diseñado sobre el modelo linealizado estabiliza la posición del levitante en presencia de la fricción estática y de la zona muerta, sin compensarlas explícitamente, en todas las pruebas realizadas dentro del intervalo de operación que permite el actuador. La segunda parte no se confirma: la alta ganancia a baja frecuencia que aporta la segunda integración no reduce el tiempo de atascamiento ni la amplitud del ciclo límite respecto a un controlador con una sola acción integral; con la misma ganancia, ambas estructuras producen el mismo ciclo, y la duración del atascamiento la fija el término integral simple.

\paragraph{Síntesis.}
Se diseñó y validó por simulación un controlador PII que estabiliza la posición del SLM ECP-730 en configuración inestable con atascamiento-deslizamiento, con error medio prácticamente nulo y sin alcanzar el límite superior del actuador, dentro del intervalo de operación que permite un voltaje máximo de $3{,}5$~V. Los resultados también delimitan el alcance de la propuesta: el desempeño en estado estacionario está dominado por el ciclo límite que imponen la fricción estática y la zona muerta, y la segunda acción integral no ofreció una mejora frente a un PI en ninguna de las variantes evaluadas. Estos resultados sientan las bases para la validación experimental y para las líneas de trabajo que se proponen a continuación.
